\scanpage{67}
\uline{Parties $i$-libres} etc

On étudie des types de diagrammes plongés dans les complexes cubiques $C$, liés directement aux parties $n$-libres \deleted{\unsure{ou $n$-bases}} (parties de cardinal $n$ de $C$ formées de sommets non liés deux à deux). Signalons \deleted{tout} de suite

\uline{Prop} \textcircled{1} Pour qu'il existe dans $C$ une partie $n$-libre, \deleted{\unsure{il}} faut s. que $1\leq n\leq6$. Les parties $n$-libres maximales \inserted{s}/sont : \deleted{celles}

a) celles de cardinal $6$, ce sont aussi des bases

b) \deleted{celles de cardinal $5$ qui sont \unsure{représentées}} \inserted{les parties} de la \deleted{forme} $E(x)\cap E(y)$, où $x,y$ sont deux sommets distincts \unsure{non liés|deux liés} de $C$. Ce sont des parties de cardinal $5$.

\textcircled{2} Soit $L_{\unsure{n|i}}$ \quad $1\leq\deleted{\unsure{n}}i\leq6$ le graphe discret \inserted{$[1,i]$} de cardinal $i$ (pas de liens). Alors \deleted{pour} $i\ne5$, l'ens $\mathrm{Pl}(L_n,C)$ des plongements de $L_n$ dans $C$ est un espace homogène sous $W=\mathrm{Aut}(C)$. Pour $i=5$, il a deux orbites \inserted{$\mathrm{Pl}(L_5,C)'$ et $\mathrm{Pl}(L_5,C)''$} pour $W$, savoir

a) les $5$-bases, i.e. parties $5$-libres qui sont contenues dans une base -- une !

b) les demi étoiles, i.e. parties de la forme $E(x)\cap E(y)$, où $x,y$ sont deux sommets distincts non liés

Les applications \deleted{Pl} (hom d'un pas en bas)

\diagram{\resizebox{\linewidth}{!}{$\begin{tikzcd}[ampersand replacement=\&,column sep=small,row sep=small]
{} \& {} \& {} \& {} \& {} \& \mathrm{Pl}''(L_5,C)\arrow[dl,"1"'] \& {}\\
\underset{\substack{\Vert\\\mathrm{Pl}(L_0,C)}}{\mathrm{pt}} \& \underset{\substack{\Vert\\D_1}}{\mathrm{Pl}(L_1,C)}\arrow[l,"27"'] \& \mathrm{Pl}(L_2,C)\arrow[l,"16"'] \& \mathrm{Pl}(L_3,C)\arrow[l,"10"'] \& \mathrm{Pl}(L_4,C)\arrow[l,"6"'] \& \mathrm{Pl}'(L_5,C)\arrow[l,"2"'] \& \mathrm{Pl}(L_6,C)\arrow[l,"1"']
\end{tikzcd}$}}

sont surjectives, et les degrés de projection sont \deleted{respectivement} \deleted{$27,16,\unsure{90|10},6$} ceux marqués (cf \unsure{dévissage|dénissage|dessin} plus bas), d'où on tire les cardinaux des $i$-bases \unsure{épinglées}

% archive sheet number: 66
